\documentclass[12pt,a4paper]{report}
\usepackage[utf8]{inputenc}
\usepackage{amsmath}
\usepackage{amsfonts}
\usepackage{amssymb}
\usepackage{amsthm}
\usepackage{hyperref}

\usepackage{multicol}
\usepackage{fancyhdr}
\usepackage[inline]{enumitem}
\usepackage{tikz}
\usepackage{tikz-cd}
\usetikzlibrary{calc}
\usetikzlibrary{shapes.geometric}
\usetikzlibrary{positioning}
\usepackage[margin=0.5in]{geometry}
\usepackage{xcolor}

\hypersetup{
    colorlinks=true,
    linkcolor=blue,
    filecolor=magenta,      
    urlcolor=cyan,
    pdftitle={Tensors},
    pdfpagemode=FullScreen,
    }

%\urlstyle{same}

\newcommand{\CLASSNAME}{Math 5050 -- Special Topics: Differential Geometry}
\newcommand{\STUDENTNAME}{Paul Carmody}
\newcommand{\ASSIGNMENT}{Section 12: Tangent Bundles }
\newcommand{\DUEDATE}{September 24, 2025}
\newcommand{\PROFESSOR}{Professor Berchenko-Kogan}
\newcommand{\SEMESTER}{Fall 2025}
\newcommand{\SCHEDULE}{TBD}
\newcommand{\ROOM}{Remote}

\newcommand{\MMN}{M_{m\times n}}
\newcommand{\FF}{\mathcal{F}}

\pagestyle{fancy}
\fancyhf{}
\chead{ \fancyplain{}{\CLASSNAME} }
%\chead{ \fancyplain{}{\STUDENTNAME} }
\rhead{\thepage}
\input{../MyMacros}

\newcommand{\RED}[1]{\textcolor{red}{#1}}
\newcommand{\BLUE}[1]{\textcolor{blue}{#1}}

\newcommand{\SUPP}{\operatorname{supp}}
\newcommand{\CLOSURE}{\operatorname{closure of}}

\begin{document}

\begin{center}
	\Large{\CLASSNAME -- \SEMESTER} \\
	\large{ w/\PROFESSOR}
\end{center}
\begin{center}
	\STUDENTNAME \\
	\ASSIGNMENT -- \DUEDATE\\
\end{center} 


\begin{enumerate}[label=12.\arabic*]

\item \textbf{Hausdorff condition on the tangent bundle}

Prove Proposition 12.4

\BLUE{Let $U, V$ be disjoint open sets in $M$. Then, $\pi^{-1}(U)$ and $\pi^{-1}(V)$ are open sets in $TM$ as $\pi$ is diffeomorphic.  Since $M$ is Hausdorff, choose $p \in M$ be such that $p \not \in U$ and $p \not \in V$.   Then $T_p(M) \in TM$ and disjoint from all other elements of $TM$ including the sections $\pi^{-1}(U)$ and $\pi^{-1}(V)$.  Thus, $TM$ is Hausdorff.
}

\item \textbf{Transition functions for the total space fo the tangent bundle}

Let $(U,\phi)= (U, x^1, \dots, x^n)$ and $(V, \psi) = (V, y^i, \dots, y^n)$ be overlapping coordinate charts on a manifold $M$.  They induce coordinate charts $(TU, \tilde{\phi})$ and $(TV, \tilde{\psi})$ on the total space $TM$ of the tangent bundle (see equation (12.1)), with transition function $\tilde{\psi}\circ\tilde{\phi}^{-1}$:
\begin{align*}
	(x^1,\dots, x^n, a^1,\dots, a^n)\mapsto(y^1,\dots, y^n, b^1, \dots, b^n).
\end{align*}\begin{enumerate}[label=(\alph*)]
	\item Compute the Jacobian matrix of the transition function $\tilde{\psi}\circ\tilde{\phi}^{-1}$ at $\phi(p)$.
	
	\item Show that the Jacobian determinant of the transition function $\tilde{\psi}\circ\tilde{\phi}^{-1}$ at $\phi(p)$ is $(\det[\partial y^i/\partial x^j])^2$.
	
	\BLUE{The Jacobian is the differential/pushfoward of the diffeomorphism $F: TU \to TV$ and let $T=\tilde{\psi}\circ\tilde{\phi}^{-1}$ and  
	\begin{align*}
		\tilde{x}^i &= \BINDEF{x^i & i \le n}{a^{i-n} & i > n} \\
		dT &= \sum_{i=1}^{2n}  \RBAR{\PART{T^j}{\tilde{x}^i}}_{\phi(p)} \\
		&= \sum_{i=1}^{2n}  \RBAR{\PART{\tilde{y}^j}{\tilde{x}^i}}_{\phi(p)} 
	\end{align*}We get a four quadrant matrix of the form
	\begin{align*}
		J(T) &= \PAREN{\begin{array}{cc}
			J(x,y) & J(a,y) \\
			J(x,b) & J(a,b) 
		\end{array}
		} \\
		\det J(T) &= J(x,y)J(a,b) - J(a,y)J(x,b)
	\end{align*}$J(a,y)$ and $J(x,b)$ are both zero because the vectors $a$ and $b$ are from different vector spaces than $x$ and $y$, respectively. And $J(a,b)$ is precisely $J(x,y)_{\phi(p)}$, thus
	\begin{align*}
		\det J(T(\phi(p))) &= J(x,y)_{\phi(p)}^2
	\end{align*}
	}
\end{enumerate}
	
	\item \textbf{Smoothness of scalar multiplication}
	
	Prove Proposition 12.9(ii).
	
	\item \textbf{Coefficients relative to a smoothframe}
	
	Let $\pi: E \to M$ be a $C^\infty$ vector bunlde and $s_1,\dots, s_r$ a $C^\infty$ frame for $E$ over an open set $U$ in $M$.  Then every $e \in \pi^{-1}(U)$ can be written uniquely as a linear combination
	\begin{align*}
		e = \sum_{j=1}^r c^j(e)s_j(p), \, p = \pi(e) \in U.
	\end{align*}Prove that $c^j: \phi^{-1} U \to \R$ is $C^\infty$ for $j=1,\dots, r$.  (\textit{Hint:} First show that the coefficients of $e$ relative to the frame $t_1, \dots, t_r$ of a trivialization are $C^\infty$.)


\end{enumerate}

\end{document}
